\documentclass[12pt]{article}
\usepackage[utf8]{inputenc}
\usepackage[T2A]{fontenc}
\usepackage[english,russian]{babel}
\usepackage{amsmath}
\usepackage{amssymb}
\usepackage[a4paper,margin=2cm]{geometry}
\usepackage{setspace}
\usepackage{hyperref}
\usepackage{mathtools}
\onehalfspacing
\sloppy

\title{Достаточное условие существования решения задачи Коши. Формулировка теоремы Пикара.}

\begin{document}

\maketitle

\section{Достаточное условие существования решения задачи Коши}
Рассмотрим уравнение
\begin{equation}\label{eq:general-ode}
    y' = f(x, y),
\end{equation}
где $f(x, y)$ определена в некоторой области $D \subset \mathbb{R}^2$.
Рассмотрим также точку $(x_0, y_0) \in D$ в которой $y(x_0) = y_0$.
Получаем задачу Коши:
\begin{equation}\label{eq:initial-condition-2}
    \begin{cases}
        y' = f(x, y), \\
        y(x_0) = y_0
    \end{cases}
\end{equation}
Рассматриваем $x\in (x_0, a]$, где $a > x_0$ - фиксированное число.
Считаем, что $f(x, y)$ непрерывна в области $D = [x_0, a] \times [y_0 - b, y_0 + b]$ для некоторого $b > 0$.
Получаем прямоугольник на плоскости $x0y$.

\subsection{Дополнительные условия}
\textbf{Условие 1.} $f(x, y)$ определена и непрерывна в области $D$.

Тогда, по теормеме Вейерштрасса, $f(x, y)$ ограничена в области $D$.
То есть существует такое число $M > 0$, что для всех $(x, y) \in D$ выполняется неравенство
\begin{equation}\label{eq:bound-f}
    |f(x, y)| \leq M \quad \forall (x, y) \in D.
\end{equation}

\textbf{Условие 2.} $f(x, y)$ удовлетворяет условию Липшица по переменной $y$ в области $D$.

То есть существует такое число $L > 0$, что для любых двух точек $(x, y_1) \in D$ и $(x, y_2) \in D$
выполняется неравенство
\begin{equation}\label{eq:lipshitz}
    |f(x, y_1) - f(x, y_2)| \leq L |y_1 - y_2|.
\end{equation}

\subsection{Теорема о существовании единственного решения задачи Коши (Теорема Пикара)}.
Пусть $f(x, y)$ определена и непрерывна в области $D$ и удовлетворяет условиям 1 и 2.
Тогда решение задачи Коши
\begin{equation}\label{eq:cauchy-problem-2}
    \begin{cases}
        y' = f(x, y), x \in (x_0, a] \\
        y(x_0) = y_0
    \end{cases}
\end{equation}
существует и единственно на отрезке
$[x_0, x_0 + h]$, где $h = \min\left(a - x_0, \frac{b}{M}\right)$.

\end{document}